\section{Introduction}
\label{sec:introduction}

GPU implementations of entropic optimal transport (EOT) have become markedly
better at executing a single Sinkhorn update. Matrix-free operators, tiled
reductions, and fused streaming kernels reduce memory traffic and avoid
materializing the transport matrix
\citep{feydy2019sinkhorn,cuturi2022ott}. Many workloads,
however, prepare a batch of couplings for predictive maintenance, scientific
state transitions, or minibatch training. Once the inner update is fast, total
cost also depends on how long each problem remains in that batch.

Different problems usually reach the stopping tolerance at different
iterations. A static batch runs until the problem requiring the most iterations
finishes: completed problems continue to occupy state, verification bandwidth,
and kernel lanes. A logical
mask freezes their state but does not shrink the tensor seen by a fused kernel.
The updates between each problem's completion and the end of the batch are
wasted OT work.

We introduce \policy{}, an execution layer that removes this padding. It starts
from standard batched execution of independent Sinkhorn problems and removes a
problem after it passes verification used by the selected solver. A
full alternating cycle makes the newly updated marginal current in exact
arithmetic, so the other marginal provides a cheap batched screen. Lanes that
pass the subsequent two-marginal verifier are removed from potentials, marginals,
support descriptors, and scratch buffers. The next backend update runs at the
width of the unfinished set.

This design separates kernel optimization from batch execution.
The inner backend determines the cost of one stabilized update; \policy{}
changes the number and width of updates executed across a batch. We implement
the execution layer with a packed Triton kernel, OTT-JAX, and PyKeOps. The packed
Triton path additionally uses the current marginal in a preliminary screen; the
OTT-JAX and PyKeOps paths remove completed problems with independent backend
mechanisms.

For problem $t$, let $n_t$ be the batched-iteration number of the first
configured verification round at which it passes the two-marginal verifier.
Fixed-width execution performs $Wn_{\max}$ problem updates, whereas dynamic
compaction performs exactly $\sum_t n_t$. This identity depends only on the
observed iteration numbers. Removable work exists whenever a batch contains
both completed and unfinished problems before an update. If
$c_k(a,\xi_k)$ denotes the measured cost of update $k$ at width $a$, the
wall-clock saving is
$\sum_k[c_k(W,\xi_k)-c_k(A_k,\xi_k)]$ minus incremental control overhead.
The completion depths and width-dependent update times therefore determine the
runtime benefit.

GPU update time can change at scheduling-wave boundaries rather than in
proportion to the number of problems. The packed Triton update launches a
width-dependent number of thread blocks; reducing $A_k$ saves time when the
grid requires fewer waves. We measure this relation with a width sweep.

Our contributions are:
\begin{enumerate}[leftmargin=*,itemsep=2pt,topsep=3pt]
  \item \textbf{Dynamic compaction during batched Sinkhorn.} Starting from
  standard fixed-width batching and the selected solver's stopping rule, we
  dynamically remove completed problems from all per-problem state. Each
  subsequent kernel executes at a width of the unfinished set.
  \item \textbf{Hardware-aware performance model.} An update-count identity
  measures removed work. A matched runtime decomposition uses measured update
  time at each active width and includes control overhead.
  \item \textbf{Component analysis.} Matched controls separate
  batched verification, Sinkhorn-specific screening, logical masking, and dynamic compaction. Grouping experiments that hold total iteration counts fixed and
  width sweeps identify when dynamic compaction reduces runtime.
  \item \textbf{Evaluation across backends and tasks.} MetroPT-3, Packer19, and
  ImageNet-32 evaluate the packed Triton path; OTT-JAX and PyKeOps test transfer
  on ImageNet-32 and A2D2. The evaluation reports runtime, residuals,
  finite-precision decisions, consumer outputs, and training quality.
\end{enumerate}
